% !TEX root = ../main.tex
\chapter{修改与提交前自查}
\label{chap:guide3}

\section{从内容关系开始安排修改顺序}

修改可以从影响范围较大的问题开始。先检查研究问题是否明确，章节是否共同回答这个问题，再检查段落之间的连接，最后检查措辞、标点与版式。若论证结构仍在变化，过早逐字调整格式可能带来重复工作。每轮修改可以列出一个明确目标，例如检查所有结论是否有对应材料，或者统一变量名称，而不是笼统地要求把全文变得更好。

图~\ref{fig:workflow} 表现“材料整理、初稿、检查、修改、编译”这五个步骤。检查发现问题后，可以回到相关材料或章节重新处理；流程中的编译负责验证当前工程能否产生排版结果，不承担学术评价。把修改原因和涉及文件记下来，有助于后续解释某段内容为什么发生变化，也便于区分本轮新出现的问题与此前已经记录的待办。

\begin{figure}[htbp]
\centering
\fbox{\parbox[c][12mm][c]{.13\linewidth}{\centering 材料\\整理}}\hfill
$\rightarrow$\hfill
\fbox{\parbox[c][12mm][c]{.13\linewidth}{\centering 初稿}}\hfill
$\rightarrow$\hfill
\fbox{\parbox[c][12mm][c]{.13\linewidth}{\centering 检查}}\hfill
$\rightarrow$\hfill
\fbox{\parbox[c][12mm][c]{.13\linewidth}{\centering 修改}}\hfill
$\rightarrow$\hfill
\fbox{\parbox[c][12mm][c]{.13\linewidth}{\centering 编译}}
\caption{从材料整理到编译的示例写作流程}
\label{fig:workflow}
\end{figure}


\section{将待确认事项保留在可检索的位置}

待办项应说明缺少什么、需要向谁核对以及影响哪些内容。只写“待完善”难以指导下一步行动；写明“核对提交要求中的摘要长度与附件清单”则更容易执行。待办既可以出现在相关段落旁，也可以汇总到交付说明中，但两处内容应保持一致。不能因为暂时无法获取答案，就用一个看似合理的数字或规则代替它。

本指南面向通用写作过程，尚未取得具体院系的最新要求和具体导师对章节安排的意见。因此保留以下两项待办：

\begin{itemize}
\item {}[TODO: 核对所在院系的最新提交要求。]
\item {}[TODO: 确认导师对章节安排的具体要求。]
\end{itemize}

完成这些待办需要读者结合自己的实际情况提供信息。公开教学示例保留它们，是为了说明缺项如何被发现和追踪，并不意味着所有论文都应使用同一份检查清单。确认后，应把依据、修改位置和结果一起记录，再运行与修改内容相关的检查。待办消失本身不能证明答案可靠，来源仍然需要审阅。

\section{分别报告内容审阅、结构检查和编译}

提交前可以从三个方面检查当前成果。内容审阅关注论点、证据和表述是否一致；结构检查关注文件路径、章节连接、标签与引用等是否完整；真实编译关注工具是否成功生成本次排版产物。三者相互补充，但不能相互替代。检查工具报告结构通过时，作者仍应阅读正文；生成 PDF 后，也应逐页观察公式、图表和段落是否发生裁切或错位。

对检查结果的记录应包含本次使用的材料版本、检查范围、实际执行的命令、发现的问题和未验证的部分。如果仍有待办，应明确说明哪些内容尚待确认。重复编译时，要确认看到的是本次生成的文件，而不是早先留下的 PDF。这样，交付者与接收者才能围绕同一组材料讨论，避免把文件存在、检查通过和正式提交混为一谈。

本指南最终交付的是可以继续编辑的示例工程。读者可用它理解章节组织与检查方法，再把这些方法应用到自己的材料。公开展示应突出原创正文、公式、表格和自查流程；真实姓名、学号以及未经确认的流程页面不应成为示例内容。具体论文的学术判断、院系规范和导师意见仍需由作者在实际写作中逐项落实。
